\section{Implementation Challenges \& Future Work}

The development of the Galois Execution Engine presented several non-trivial implementation and performance challenges, particularly concerning the $\texttt{Key-Scan}$ operator and database resource management.

\subsection{Performance Deficiencies of the $\texttt{Key-Scan}$ Operator}

The most significant performance challenge was encountered during the development and testing of the $\texttt{Key-Scan}$ physical operator, where initial evaluations demonstrated \textbf{poor retrieval effectiveness} under the $\texttt{GALOIS F}$ (Fact-checking) and $\texttt{GALOIS WO}$ (Write-out) logical plans.

\subsubsection{The Locality-of-Reference Problem}
The poor performance stems from a fundamental design aspect of the $\texttt{Key-Scan}$'s two-phase strategy: the iterative Phase 1 is designed solely to retrieve unique primary key values ($\mathcal{K}_{all}$) without applying the initial query's non-key predicates.
\\\\
For instance, consider querying the $\texttt{venezuela\_presidents}$ table with $\texttt{president\_name}$ as the primary key. In Phase 1, the LLM is prompted to retrieve $\textit{all}$ keys globally available in the training data (e.g., all world presidents). The actual filter condition (e.g., $\texttt{country} = \text{'Venezuela'}$) is applied only in the post-retrieval Phase 2 (the Tail Request). If the LLM's initial response set $\mathcal{K}_{all}$ lacks sufficient locality (i.e., fails to yield keys related to 'Venezuela'), the final result set, post-query filtering, will be empty ($\mathcal{R}_{final} = \emptyset$). This retrieval inefficiency was common across several datasets used in the $\texttt{GALOIS F}$ and $\texttt{GALOIS WO}$ contexts, necessitating future work focused on improving the $\texttt{Key-Scan}$'s key retrieval effectiveness.

\subsection{Database Resource Management for $\texttt{DuckDB}$}

A critical architectural challenge involved managing the concurrent access and connection lifecycle for the $\texttt{DuckDB}$ backend. The execution pipeline requires numerous connection instantiations across different classes for tasks such as retrieving column data types, inspecting the database schema, and executing final queries.
\\\\
To address the potential for file locks and resource contention on the underlying $\texttt{.duckdb}$ instance, the resource management for the $\texttt{DuckDB}$ backend adheres to a strict \textbf{acquisition-release pattern}:
\begin{itemize}
    \item \textbf{Scoped Instantiation:} Database connections are instantiated solely within the confined scope of the active schema retrieval or query execution task.
    \item \textbf{Guaranteed Release:} Connections are immediately released within $\texttt{finally}$ blocks upon completion.
\end{itemize}
This explicit lifecycle management guarantees that the transition between the schema inspection phase and the subsequent in-memory virtualization phase proceeds without contention, maintaining system stability.

\subsection{Addressing Sub-Optimal Prompt Structure}

The core functionality of the $\texttt{Table-Scan}$ and $\texttt{Key-Scan}$ iterative models has been established; however, current system performance indicates that the \textbf{structure and semantic design of the prompts} generated by the existing Galois framework are sub-optimal. These prompts fail to consistently elicit high-fidelity, complete structured output from the LLMs.
\\\\
This observation points to a critical area for future work: systematically improving the retrieval effectiveness by focusing on prompt engineering. Future efforts will concentrate on refining the instructions, contextualization ($\mathcal{H}$ and $\mathcal{P}_{pushdown}$ injection), and required output formats (e.g., explicit JSON schemas) to ensure the LLM provides more precise and comprehensive results. The hypothesis is that \textbf{optimized prompt semantics} will directly lead to a reduction in the number of required iterations, lower token consumption, and ultimately, higher data coverage.
